\begin{enumerate}
    \item \textbf{Gate (RQ1):} an approximate RBF SVM can serve as a 98\%-recall gate at full CICIDS2017 scale. Its false-positive rate depends strongly on hyperparameters: validation-based tuning reduced it from \SvmValFprBefore\% to \SvmValFprAfter\% (test: \EoneGateFpr\%). It remains the main source of false alarms and missed web attacks.
    \item \textbf{Stacking (RQ2):} family attribution is above 99.5\% accurate for all models, but out-of-fold stacking did not improve macro-F1 over its best base learner, the Random Forest (\EoneStackMacroFone{} vs.\ \EoneRfMacroFone), because of very rare families.
    \item \textbf{Autoencoder (RQ3):} the benign-only autoencoder is weak alone but complementary. It turns most rare-family flows missed by the gate (Heartbleed, Infiltration) into \texttt{anomaly\_review} alerts.
    \item \textbf{Baselines (RQ4):} single tree ensembles (XGBoost: \EoneXgbRecall\% recall at \EoneXgbFpr\% FPR) beat the fused engine on binary detection in the closed-set setting. The engine's advantages are explainable, staged evidence and family attribution, at \EoneMsPerFlow{} ms per flow.
    \item \textbf{Generalisation (RQ5):} on an unseen capture day with an unseen attack family, every model degrades. Tree ensembles miss over a third of the DDoS flows, the gate's validation-calibrated threshold flags most benign traffic, and benign traffic itself shifts. Results from random splits of CICIDS2017 (where \EoneOverlap\% of test flows have an exact copy in training) substantially overstate real-world performance.
\end{enumerate}
